\section{The venue says what it is doing}
\label{sec:ha:technical_event}

The venue runs one trading day at a time. During the day it passes through five phases, and it
tells its components and its members which phase it is in and whether trading is halted. This
section is in the high availability chapter because \reqref{R-0126} requires the venue to report
when it cannot do a component's work, and the venue cannot report that until it can announce
anything about itself at all.

The design these requirements come from is in \texttt{docs/venue/trading\_phases.md}, which sets
out the phases and the record that announces them, and \texttt{docs/venue/trading\_day\_plan.md},
which sets out the whole day.

\subsection{The phases and the record that announces them}
\label{sec:ha:technical_event:phases}

The venue is always in exactly one of five phases, which follow one another in this order:

\begin{description}
  \item[Start of day.] The venue is preparing for the day and is not trading.
  \item[Pre-open.] Members may enter and cancel orders. A venue that matches orders would match
    nothing yet, and would run an opening auction when it moves to the next phase.
  \item[Open.] Members may enter and cancel orders, and the venue trades.
  \item[Closed.] Trading has finished for the day. On entering this phase the venue cancels every
    order still open.
  \item[End of day.] Nothing further is sequenced, and the day's record is completed.
\end{description}

A halt is not a phase. Whether trading is halted is a flag held alongside the phase, so that
lifting a halt returns the venue to the phase it is in, with nothing remembered about the phase it
was in before.

The leading sequencer writes a \texttt{TechnicalEvent} command into the WAL for each change of
phase and for each halt and lift. The record states the phase the venue is now in, whether trading
is halted, and the trading day. A component that applies one record therefore knows the venue's
whole state, even if it missed every record before it.

The table gives what each phase allows.

\begin{center}
\small
\begin{tabular}{lccccc}
  & Start of day & Pre-open & Open & Closed & End of day \\
  \hline
  A member may log on                & yes & yes & yes & yes & yes \\
  A member may enter an order        & no  & yes & yes & no  & no  \\
  A member may cancel an order       & no  & yes & yes & no  & no  \\
  The matching engine acts on orders & no  & yes & yes & no  & no  \\
\end{tabular}
\end{center}

This venue's matching engine does no matching, so it treats pre-open and open alike. Members may
log on in every phase, so that a member is told what the venue is doing and why its orders would be
refused.

\subsection{The trading day}
\label{sec:ha:technical_event:trading_day}

\begin{requirement}{R-0164}{The trading day is given to the run, and no component computes it}
  The procedure that starts the venue shall give it the trading day. The first record in the
  day's WAL shall be the \texttt{TechnicalEvent} that enters start of day for that trading day,
  and every later record shall carry the same trading day. No component shall compute the
  trading day or the phase from its own clock: each shall hold the values the venue declared.
  \rationale{The trading day is a business date and need not match the calendar date on any
  machine. A session can open the evening before, and a machine in another time zone can turn
  over to the next date while the venue is in the middle of its day. Two components computing
  the trading day from their clocks near a boundary would disagree, each correct by its own
  clock, and nothing would detect it. A replay of an earlier day is run on a later date, and must
  produce the same records.}
  \uncovered{no scenario starts the venue with a given trading day}
\end{requirement}

\begin{gap}{BUG-0137}
  The leading sequencer takes the trading day from its clock: the date, in UTC, of the first
  record in the day's WAL. There is no start-of-day \texttt{TechnicalEvent}, and no procedure gives
  the venue its trading day.
\end{gap}

\begin{requirement}{R-0165}{A record for an older trading day is refused}
  A component shall refuse a \texttt{TechnicalEvent} whose trading day is older than the trading
  day it holds, and shall report that it refused it.
  \rationale{The trading day plays the part here that the epoch plays in leadership. A halt
  declared yesterday must not stop trading today, and a record from a WAL left over from an
  earlier day must not change the state of the present one.}
  \uncovered{the matching engine refuses such a record, and only a unit test checks it}
\end{requirement}

\subsection{How the venue moves from one phase to the next}
\label{sec:ha:technical_event:moving}

\begin{requirement}{R-0166}{The phases follow one another in order, each entered once}
  The venue shall enter the phases in the order start of day, pre-open, open, closed, end of
  day, and shall enter each at most once in a trading day. The leading sequencer shall write the
  \texttt{TechnicalEvent} for a phase only while the venue is in the phase before it, and only
  when that phase's conditions hold. Entering pre-open requires the matching engine to have
  applied every record that must be in place before members enter orders. Entering end of day
  requires every report produced by the close to be in the WAL.
  \rationale{A phase that could be entered twice, or skipped, would leave the components to work
  out which of two contradictory states is current. Writing a change only from the phase before
  it is also what makes a repeated request, and a change of leader, harmless: the second attempt
  finds the venue already moved on and writes nothing.}
  \uncovered{no phase other than open is written yet}
\end{requirement}

\begin{requirement}{R-0167}{The venue follows a schedule recorded in its WAL}
  The day's schedule shall be given to the venue with the request that starts the day, as the
  earliest time at which each phase may be entered. The leading sequencer shall write the
  schedule into the WAL. When the time for the next phase has come and that phase's conditions
  hold, the leading sequencer shall write the \texttt{TechnicalEvent} that enters it. No other
  component shall change the phase because of its own clock.
  \rationale{A clock on one machine is not a valid cause of a change of state in a replicated
  venue, because the other machines and a later replay do not share it. Recording the schedule in
  the WAL means every sequencer holds the same times, a newly promoted leader knows when the next
  change is due, and the WAL shows why each change was made. The leading sequencer already
  carries out a member's cancel-on-disconnect grace period in the same way: the deadline is in the
  WAL, and the leader writes the cancel when the deadline passes. Aeron Cluster writes the expiry
  of a timer into its replicated log for the same reason.}
  \uncovered{no schedule exists yet}
\end{requirement}

\begin{requirement}{R-0168}{An operator can move the venue on, and change the schedule}
  An operator shall be able to ask for the next phase to be entered at once, and to replace the
  times remaining in the day's schedule. The leading sequencer shall answer every such request
  from its current state, saying whether it wrote the change, whether the venue is already in
  or past the phase asked for, or why it refused. A sequencer that does not lead shall say so.
  \rationale{Real venues publish a schedule and override it when they need to: the London Stock
  Exchange describes market control actions that override the normal schedule. An early close is
  the common case. Answering from the current state makes a repeated request harmless, so an
  operator whose answer was lost asks again and learns where the venue now is.}
  \uncovered{no request to change the phase exists yet}
\end{requirement}

\begin{requirement}{R-0169}{The phase moves on while halted, except when the orders cannot be accounted for}
  The venue shall move from one phase to the next while trading is halted, except while it is
  halted because a matching engine cannot account for the orders the venue holds
  (\reqref{R-0123}). In that case the phase shall not move on until a person has resolved the
  halt.
  \rationale{A halt at 16:20 must not stop the close at 16:30 from cancelling the open orders,
  or the day from ending. When an engine cannot account for the orders, nobody knows which orders
  are open, so the close cannot cancel every open order and tell each member, and a person has to
  resolve what the venue holds first.}
  \uncovered{no phase other than open is written yet}
\end{requirement}

\subsection{What each phase allows}
\label{sec:ha:technical_event:allows}

\begin{requirement}{R-0170}{The sequencer refuses a command the phase does not allow}
  The leading sequencer shall refuse every command from a member that the current phase does not
  allow, as the table in \cref{sec:ha:technical_event:phases} gives, and shall say which phase
  caused the refusal.
  \rationale{A gateway acts on what the sequencer last told it, and that can be out of date by
  the time a command arrives. The sequencer holds the phase from its WAL, so its refusal is the
  one no out-of-date view can get past.}
  \uncovered{the sequencer refuses no command because of the venue's state}
\end{requirement}

\begin{requirement}{R-0171}{A gateway refuses a command the phase does not allow, and says why}
  The leading sequencer shall tell every gateway the trading day, the phase and whether trading is
  halted, when the gateway connects and whenever any of them changes. A gateway shall refuse a
  command the phase does not allow before passing it to the sequencer, and shall tell the member
  the reason. Both gateways shall refuse the same commands for the same reasons
  (\reqref{R-0152}).
  \rationale{A command refused at the gateway costs the sequencer nothing. A member must be able
  to tell a command refused because the market is closed from one refused because trading is
  halted, because the first is expected every evening and the second needs attention.}
  \uncovered{the gateways learn whether trading is halted, not the phase}
\end{requirement}

\begin{requirement}{R-0172}{The close cancels every open order, and each member is told}
  The \texttt{TechnicalEvent} that enters the closed phase shall be the command that cancels every
  order still open. Each matching engine, applying it at its position in the WAL, shall cancel
  every order on its book and produce an execution report for each, to the member that placed
  it.
  \rationale{An order that outlives the day would be open overnight with nobody watching it, in a
  venue that starts the next day with an empty book. Because the cancellation is one command at
  one position, both matching engines cancel the same orders, and a replay cancels them again in
  the same way. Every member's view of its orders comes to match the venue's, for the reason
  \reqref{R-0024} gives.}
  \uncovered{the close does not exist yet}
\end{requirement}

\begin{gap}{BUG-0132}
  Orders still open at the end of the trading day are never cancelled, and their members are
  never told.
\end{gap}

\begin{requirement}{R-0173}{No order record follows end of day}
  No record of an order, a cancel or an execution report shall follow the \texttt{TechnicalEvent}
  that enters end of day in the WAL.
  \rationale{The order activity recorder completes the day when every event up to the end-of-day
  event is confirmed by the external messaging system. That means the whole day is recorded only
  if nothing that concerns an order comes after it. The requirements above ensure it: no command
  from a member is accepted after the close, and end of day is not entered until every report the
  close produced is in the WAL.}
  \uncovered{the end-of-day phase does not exist yet}
\end{requirement}

\subsection{Who is told}
\label{sec:ha:technical_event:told}

There are two audiences. Components that act on the phase, the sequencer and the gateways, take it
on the control connection between them. Components that only need to know it take it by publish and
subscribe.

\begin{requirement}{R-0174}{Every logged-on member is told of each change}
  A gateway shall tell every logged-on member the phase and whether trading is halted each time
  either changes, and shall tell each member the current state as it logs on. A member shall be
  able to ask for the current state at any time, and the gateway shall answer from the state it
  holds. The FIX gateway shall use \texttt{TradingSessionStatus} (35=h) and
  \texttt{TradingSessionStatusRequest} (35=g); the binary gateway shall use equivalent messages.
  \rationale{A member that is connected and not trading sends no order, so the refusal of an
  order tells it nothing. A member that logs on during a halt, or before the open, would otherwise
  learn the venue's state only by sending an order and having it refused. Every member is told
  every change, with no list of which members asked, because every member wants to know when the
  market closes.}
  \uncovered{neither gateway sends or understands 35=h or a binary equivalent}
\end{requirement}

\begin{gap}{BUG-0065}
  A member learns of a halt only from the reply to an order, and is never told the phase.
\end{gap}

\begin{requirement}{R-0175}{Components that only need to know take the changes from a topic}
  The matching engine publisher shall publish every \texttt{TechnicalEvent} on a topic of its own,
  carrying the sequence number the sequencer gave it. A component that has just started shall be
  able to read the topic from the beginning of the day and so learn the venue's state without
  waiting for the next change.
  \rationale{The order activity recorder and a market data feed act on no member's command but
  need each change in its place among the orders and execution reports. The sequence numbers give
  that place across topics. A separate topic keeps each topic to one kind of record, and adds no
  load to the sequencers, which a component following their WAL itself would.}
  \uncovered{the matching engine publisher publishes no technical events}
\end{requirement}

\subsection{What is not specified here}
\label{sec:ha:technical_event:not_here}

A phase or a halt for one instrument is a different message from the venue's, as FIX separates
\texttt{SecurityStatus} (35=f) from \texttt{TradingSessionStatus} (35=h). This section specifies
the venue's phase only, and a status for one instrument can be added without changing it.

How the order activity recorder's report that the day is completely recorded reaches the end-of-day
procedure is specified with the recorder.
